%% ma: L11-B16-0086
%% lop: 11
%% chuong: 5
%% bai: 16
%% dang: 11.16.1
%% loai: TL
%% muc: VD
%% dap_an: "-"
%% nguon: "Ảnh giáo viên gửi 25/11/2026: Giới hạn hàm số dạng 0/0 (3 bài, có ĐS), Bài 2 câu 4"
%% kien_thuc: "Lớp 11 Bài 16 (giới hạn của hàm số tại một điểm); chia đa thức theo sơ đồ Horner, phân tích đa thức thành nhân tử"
%% trang_thai: chinh-thuc
\begin{ex}
Tính giới hạn $A=\lim\limits_{x\to1}\dfrac{x^4-x^3-x+1}{x^3-5x^2+7x-3}$.
\loigiai{$x^4-x^3-x+1=(x-1)^2(x^2+x+1)$, $x^3-5x^2+7x-3=(x-1)^2(x-3)$. Biểu thức bằng $\dfrac{x^2+x+1}{x-3}$. Vậy $A=-\dfrac32$.}
\end{ex}
